\chapter[My second results chapter]{Long version of my second results chapter title}
    \label{cha:second-results}

\epigraph{%
    Dans la vie, rien n'est à craindre,\\ tout est à comprendre.\footnotemark}
    {Marie Skłodowska-Curie}

\footnotetext{\emph{%
    Nothing in life is to be feared, it is only to be understood.}}

% Remarks left here as an example
\begin{remarks}
    Example remarks: Barring minor adaptations, this chapter has been submitted to \textsc{arXiv} as is~\ccite{wardaGetting2025} and is currently under review for publication in the \textit{Journal of Chemical Theory and Computation}.
    %
    Kinga Warda carried out the majority of the \ac{DFT} calculations, while I provided guidance on the choice of the Hubbard manifold, performed the symmetry analysis of the eigenstates, and ran the HSE06 calculations. 
    %
    I also identified the spurious Hubbard force contributions described in \cref{sec:second-spurious}, created three figures (using a plotting utility developed by Iurii Timrov), and authored the majority of the manuscript. 
    %
    Iurii Timrov, Lucio Colombi Ciacchi and Piotr Kowalski supervised the project and assisted in the preparation of the manuscript.
\end{remarks}

% The following space can be used as a "section abstract" where you summarize the work and its main findings (here some parts left as an example)
%
\noindent \Ac{DFT} Lorem ipsum dolor sit amet, consetetur sadipscing elitr, sed diam nonumy eirmod tempor invidunt ut labore et dolore magna aliquyam erat, sed diam voluptua.
    %
    At vero eos et accusam et justo duo dolores et ea rebum.
    %
    Stet clita kasd gubergren, no sea takimata sanctus est Lorem ipsum dolor sit amet.
    %
    Lorem ipsum dolor sit amet, consetetur sadipscing elitr, sed diam nonumy eirmod tempor invidunt ut labore et dolore magna aliquyam erat, sed diam voluptua.
    %
    At vero eos et accusam et justo duo dolores et ea rebum.
    %
    Stet clita kasd gubergren, no sea takimata sanctus est Lorem ipsum dolor sit amet. \medskip

%%%%%%%%%%%%%%%%%%%%%%%%%%%%%%%%%%%%%%%%%%%%%%%%%%%%%%%%%%%%%%%%%%%%%%

\section{Initial tests}
    \label{sec:second-tests}
    %
    \subsection{Great resistivity}
        \label{subsec:second-resistivity}
    %
    \subsection{Not-so-great conductivity}
        \label{subsec:second-conductivity}

\section{Spurious contributions}
    \label{sec:second-spurious}